\documentclass[11pt]{amsart}
\usepackage{amssymb}
\usepackage[margin=1.15in]{geometry}
\usepackage[hidelinks]{hyperref}
\newtheorem{theorem}{Theorem}
\newtheorem{proposition}[theorem]{Proposition}
\newtheorem{lemma}[theorem]{Lemma}
\newtheorem{corollary}[theorem]{Corollary}
\theoremstyle{definition}
\newtheorem{definition}[theorem]{Definition}
\theoremstyle{remark}
\newtheorem{remark}[theorem]{Remark}
\newcommand{\code}[1]{\texttt{#1}}
% stamp on every page: file name and date/time of compilation
\newcount\stamphour \newcount\stampmin
\stamphour=\time \divide\stamphour by 60
\stampmin=\time \multiply\stamphour by 60 \advance\stampmin by -\stamphour \divide\stamphour by 60
\edef\stamptime{\the\year-\ifnum\month<10 0\fi\the\month-\ifnum\day<10 0\fi\the\day\space
 \ifnum\stamphour<10 0\fi\the\stamphour:\ifnum\stampmin<10 0\fi\the\stampmin}
\AddToHook{shipout/foreground}{\put(0.5\paperwidth,-\dimexpr\paperheight-0.45in\relax)%
 {\makebox[0pt]{\scriptsize\ttfamily\jobname.pdf\quad\stamptime}}}
\newcommand{\Gr}{\mathrm{Gr}}
\newcommand{\Hom}{\mathrm{Hom}}
\newcommand{\Det}{\mathrm{Det}}
\newcommand{\Bases}{\mathrm{Bases}}
\newcommand{\RegGcd}{\mathrm{RegGcd}}
\newcommand{\Ex}{{\textstyle\bigwedge}}
\newcommand{\trunc}[1]{\lVert #1\rVert}
\title[Finite free resolutions without bases]{Finite free resolutions without bases:\\ the structure theorem of Buchsbaum and Eisenbud and the gcd}
\date{October 26, 2026}
\begin{document}
\begin{abstract}
Coquand and Quitt\'e~\cite{CQ} gave elementary constructive proofs of
results of Buchsbaum and Eisenbud~\cite{BE73,BE74} on finite free
resolutions: a resolution of an ideal $(a_1,\dots,a_N)$ yields a greatest
common divisor of the $a_j$, computed as the Cayley determinant of the
resolution.  The statements are about matrices, and the Cayley determinant
depends on the bases, up to a unit.  We show what this ambiguity is: the
canonical object is a line with an injective map $a_1:\Det(F)\to R$, and
choosing a basis of the line produces an ordinary greatest common divisor.
The regular greatest common divisors of the $a_j$ are exactly the values of
$a_1$ at the bases of $\Det(F)$, and the pair $(\Det(F),a_1)$ depends only
on the ideal, up to a unique isomorphism.  To obtain this we restate the
results for finite free modules, in univalent foundations, following the
form of the first structure theorem that Buchsbaum and Eisenbud give for
projective modules \cite[Theorem~3.1 and the Remark after it]{BE74}: the
multipliers $a_k$ are maps between lines defined without tensor products
or duals, and the type of such structures on a complex is contractible.
The same method gives Wiebe's theorem~\cite{CT} without bases.  All
results are checked in the proof assistant Unit.
\end{abstract}
\maketitle

\section{Introduction}

Let $R$ be a commutative ring and $0\to F_n\xrightarrow{f_n}\cdots
\xrightarrow{f_1}F_0$ an exact complex of finite free modules.
Buchsbaum and Eisenbud~\cite[Theorem~3.1]{BE74} show that each
$\Ex^{r_k}f_k$ factors through a map $a_k$ into $\Ex^{r_k}F_{k-1}$
from a module of rank one, and that these factorisations are linked from
one step to the next.  In their remark after Theorem~3.1 they explain
that for projective modules, which need not be oriented, the
identifications $\Ex^{r_{k+1}}F_k\cong\Ex^{r_k}F_k^*$ have to be replaced
by natural isomorphisms involving lines
$R_k=\Ex^{s_k}F_k\otimes R_{k+1}^*$; the theorem then holds with maps
$a_k:R_k\to\Ex^{r_k}F_{k-1}$.

Coquand and Quitt\'e~\cite{CQ} prove these results constructively, with
the notion of true grade instead of depth and without prime ideals, in
coordinates: $F_k=R^{p_k}$, $f_k$ a matrix $A_k$.  Their Lemma~8.4 and
Corollary~8.5 write $\Ex^{r_k}(A_k)=\beta_k\beta_{k+1}^{*T}$ with
$\beta_k$ in $\Ex^{r_k}(F_{k-1})$; the last vector $\beta_1$ is the
Cayley determinant, a greatest common divisor of the $r_1$-minors of
$A_1$ (Theorem~8.7), and a finite free resolution of
$(a_1,\dots,a_N)$ gives a greatest common divisor of the $a_j$
(Theorem~8.9).  Two features are tied to the coordinates: the dual
$\beta^*$ is the Hodge dual, so the factorisation carries a sign
$(-1)^{pq}$ depending on conventions, and the Cayley determinant depends
on the bases up to a unit.  The book of Coquand and Lombardi~\cite{CL}
gives a complete constructive treatment of Northcott's theory~\cite{N},
including these results.

This paper gives the coordinate-free form, and with it an explanation of
the ambiguity ``up to a unit''.  For a resolution $F$ of a nonzero ideal
$(a_1,\dots,a_N)$ the canonical object is not an element of $R$ but a line
$\Det(F)$ together with an injective linear map $a_1:\Det(F)\to R$ whose
image is the gcd ideal.  A basis $\varphi$ of the line gives the element
$a_1(\varphi)$, and
\[\Bases(\Det F)\;\simeq\;\RegGcd(a_1,\dots,a_N),\qquad\varphi\mapsto
a_1(\varphi)\]
(Theorem~\ref{thm:gcd}): the regular greatest common divisors are exactly
the values of $a_1$ at the bases, and the unit relating two of them is the
unit relating two bases of a line.  Moreover the pair $(\Det(F),a_1)$
depends only on the ideal: for two resolutions $F,G$ there is a unique
isomorphism $e:\Det(F)\cong\Det(G)$ with $a_Ge=a_F$
(Theorem~\ref{thm:compare}).

The setting is univalent foundations~\cite{HoTT}: a finite free module is
a module that is \emph{merely} isomorphic to some $R^n$, with no basis
chosen.  The main steps are the following.
\begin{enumerate}
\item The lines of Buchsbaum and Eisenbud can be written with $\Hom$
  alone: $R_{n+1}=R$ and $R_k=\Hom(R_{k+1},\Ex^{s_k}F_k)$.  The
  condition on the $a_k$ is
  \[\Ex^{r_k}f_k(x)=a_k\bigl(\rho\mapsto x\wedge a_{k+1}(\rho)\bigr),\]
  which uses only the exterior product.  Under the hypotheses of
  \cite[Corollary~8.5]{CQ}, stated with ideals of minors defined without
  bases, the type of families $(a_k)$ satisfying it is contractible
  (Theorem~\ref{thm:structure}).  In coordinates the condition is
  exactly the factorisation of \cite[Lemma~8.4]{CQ} with no sign
  (Proposition~\ref{prop:dictionary}).
\item $\Det(F):=R_1$ is a line and $a_1$ is canonical.  For a resolution
  of a nonzero ideal $(a_1,\dots,a_N)$ the last rank is $1$, so
  $a_1:\Det(F)\to R$, and $a_1$ is injective.  The regular greatest
  common divisors of the $a_j$ are exactly the elements $a_1(\varphi)$,
  $\varphi$ a basis of the line $\Det(F)$, and they form an
  $R^\times$-torsor (Theorem~\ref{thm:gcd}).  The Cayley determinant is
  $a_1$ at the basis element determined by the chosen bases.
\item Two resolutions of the same nonzero ideal, of any lengths, have
  canonically isomorphic determinant lines, compatibly with $a_1$
  (Theorem~\ref{thm:compare}); an isomorphism of resolutions is an equality
  of resolutions (structure identity principle), and the induced
  isomorphism is the canonical one (Theorem~\ref{thm:natural}).  A
  resolution also decides whether the ideal is zero
  (Proposition~\ref{prop:decide}).
\item The same method gives Wiebe's theorem for linear forms on a finite
  free module, with the determinant of an endomorphism defined on the
  top exterior power (Theorem~\ref{thm:wiebe}).
\end{enumerate}

The method is uniform: uniqueness lets coordinate constructions descend
to canonical objects, and univalence makes their invariance automatic.  An
object attached to finite free modules is defined by a property that is a
proposition, or a type that is contractible; it is then an element by
unique choice.  To prove the property it suffices to prove it after
choosing bases, and there the matrix arguments of~\cite{CQ} apply.  No
construction depends on a choice of bases, and invariance under
isomorphisms is transport.

\section{Setting}

We work in univalent type theory~\cite{HoTT}.  A type is a
\emph{proposition} if any two of its elements are equal; $\trunc{A}$ is
the propositional truncation, and ``merely $A$'' means $\trunc{A}$.  A
type is \emph{contractible} if it has an element equal to every element.
\emph{Unique choice}: if $P$ is a proposition, $\trunc{P}\to P$; in
particular an element of a contractible type, or of a proposition that is
merely inhabited, is obtained without choice.  Rings are commutative with
a set as carrier; we do not assume $0\neq1$.

\paragraph{Modules and structure identity.}  For modules $A,B$ the
identity type $A=B$ is equivalent to the type of linear bijections
$A\cong B$; equivalently the type $\sum_B(A\cong B)$ is contractible.
(The type $\sum_B(A=B)$ is contractible for every type; the content of the
statement is that the same holds with $\cong$.)  This is the structure
identity principle, proved once for a notion of
structure built from operations, actions and points with axioms that form
a proposition~\cite{E}, via the fundamental theorem of identity
types~\cite[Theorem~11.2.2]{R}; modules and commutative rings are
instances, and the only algebraic input is that their axioms are
propositions (zero and negation are unique given associativity and
commutativity).  So
to prove a statement about all pairs $(B,\varphi:A\cong B)$ it suffices
to prove it for $(A,\mathrm{id})$ (induction on isomorphisms).  The same
holds for families: $\sum_{G'}\prod_i(G_i\cong G'_i)$ is a product of
contractible types, hence contractible.  For commutative rings,
$\sum_{S'}(S\cong S')$ is contractible as well.

A module $M$ is \emph{finite free of rank $n$} if merely $M\cong R^n$;
if $1\neq0$ the rank is unique (an injective linear map $R^{n+1+j}\to
R^n$ forces $1=0$, by McCoy's theorem).  A \emph{line} is a module
merely isomorphic to $R$.  For modules $K,M$, $\Hom(K,M)$ is a module;
if $K$ and $M$ are lines, so is $\Hom(K,M)$.  An element $\varphi$ of a
module $K$ is a \emph{basis} if $c\mapsto c\varphi$ is bijective, i.e.\
$c\varphi=c'\varphi$ implies $c=c'$ and every element is $c\varphi$ for
some $c$.  The bases of a line form an $R^\times$-torsor.

\paragraph{Exterior powers.}  For a finite free $M$, the type of exterior
algebras on $M$ (an algebra with a linear map $M\to E$, $v^2=0$,
universal for such maps) is contractible: on $R^n$ the exterior algebra
is constructed explicitly, and contractibility transports along $M\cong
R^n$.  We write $\Ex M$ for its centre.  An element of $\Ex M$ has degree
$k$ if its coordinates have degree $k$ in every basis; these elements form
$\Ex^kM$, and a linear map $f:M\to N$ induces $\Ex f$ and $\Ex^kf$.  If
$M$ has rank $n$, $\Ex^nM$ is a line, and in a basis $e$ the top
coordinate is an isomorphism $\Ex^nM\cong R$.

\paragraph{Grade.}  For a family $(b_i)_{i\in I}$ of elements of $R$,
\emph{regular} means $xb_i=0$ for all $i$ implies $x=0$; $\Gr(b)\geq2$
means regular and every family $(y_i)$ with $b_iy_j=b_jy_i$ is merely of
the form $y_i=tb_i$.  This is the case $k=2$ of the grade
of~\cite[Section~6]{CQ}; the true grade $\Gr(b)\geq k$ of a finite
family is defined with regular sequences after adding indeterminates.

\section{Ideals of minors and ranks without bases}

\begin{definition}
Let $f:M\to N$ be linear between finite free modules and $k\geq0$.  The
family $D_k(f)$ consists of the elements
$\lambda(f(v_1)\wedge\cdots\wedge f(v_k))$, for $v_1,\dots,v_k\in M$ and
$\lambda:\Ex N\to R$ linear.
\end{definition}

\begin{proposition}\label{prop:dk}
For bases of $M$ and $N$, with $A$ the matrix of $f$, every $k$-minor of
$A$ is an element of $D_k(f)$, and every element of $D_k(f)$ is in the
ideal generated by the $k$-minors.  Hence ``$D_k(f)$ is regular'' and
``$\Gr(D_k(f))\geq2$'' hold if and only if they hold for the $k$-minors
of $A$, and do not depend on the bases.
\end{proposition}

\begin{proof}
A minor is the coefficient at a $k$-subset of rows of the wedge of $k$
columns; take the $v_t$ to be the basis vectors of these columns and
$\lambda$ the corresponding coordinate.  Conversely the coordinates of
$f(v_1)\wedge\dots\wedge f(v_k)$ are in the ideal of $k$-minors, and a
linear form on $\Ex N$ is a combination of coordinates.  Regularity and
$\Gr\geq2$ pass to a family $b$ from a family $a$ when every $a_i$ is in
the ideal of $b$, in the sense used by the definition of $\Gr\geq2$
(using unique choice to build the factor $t$).
\end{proof}

\begin{definition}
$f$ is \emph{stable of rank $r$} if $\Ex^{r+1}f=0$ (every
$f(v_1)\wedge\dots\wedge f(v_{r+1})$ is $0$) and $D_r(f)$ is regular.
\end{definition}

Both conditions are propositions.  If $f$ is stable of ranks $r<r'$ then
$\Ex^{r'}f=0$, so $D_{r'}(f)=0$ is regular and $1=0$.  Thus the rank is
unique when $1\neq0$.  In a basis, stability of rank $r$ is the notion
of~\cite[Section~2]{CQ} for the matrix.

\section{The structure theorem}

Let $0\to F_n\xrightarrow{f_n}F_{n-1}\to\cdots\xrightarrow{f_1}F_0$ be a
complex of finite free modules, $f_kf_{k+1}=0$, with numbers $r_k\geq0$
($1\le k\le n+1$, $r_{n+1}=0$) such that $F_k$ has rank $s_k=r_k+r_{k+1}$
for $1\le k\le n$.

\begin{definition}\label{def:structure}
The lines are $R_{n+1}=R$ and $R_k=\Hom(R_{k+1},\Ex^{s_k}F_k)$.  A
\emph{structure} on the complex is a family of linear maps
$a_k:R_k\to\Ex^{r_k}F_{k-1}$, $1\leq k\leq n+1$, with
$a_{n+1}=\mathrm{id}_R$ ($\Ex^0F_n=R$), such that for $1\leq k\leq n$
\begin{equation}\label{eq:star}
\Ex^{r_k}f_k(x)=a_k\bigl(\rho\mapsto x\wedge a_{k+1}(\rho)\bigr)
\qquad(x\in\Ex^{r_k}F_k),\tag{$*$}
\end{equation}
and $a_k$ is regular for $2\le k\le n+1$: $c\,a_k(\varphi)=0$ for all
$\varphi$ implies $c=0$.
\end{definition}

Here $x\wedge a_{k+1}(\rho)\in\Ex^{s_k}F_k$, so
$\rho\mapsto x\wedge a_{k+1}(\rho)$ is an element of $R_k$.  Since
$R_{k+1}$ and $\Ex^{s_k}F_k$ are lines, $R_k$ is a line.  For a line
$L$ and a line $K$, $L\otimes K^*=\Hom(K,L)$ canonically, and through this
identification Definition~\ref{def:structure} is the form of
\cite[Theorem~3.1(a)]{BE74} given in the Remark after it, for modules
that are not oriented: the natural isomorphism used there becomes the
exterior product.  For $n$ we get $R_n\cong\Ex^{s_n}F_n$ and
$a_n=\Ex^{r_n}f_n$, as in~\cite{BE74}.

\begin{theorem}\label{thm:structure}
Assume $\Gr(D_{r_k}(f_k))\geq2$ for $2\leq k\leq n$ and $D_{r_1}(f_1)$
regular.  Then the type of structures on the complex is contractible.
\end{theorem}

Only these hypotheses are used: that $F$ is a complex, the rank equations
$s_k=r_k+r_{k+1}$, and the two grade conditions.  Neither exactness nor
stability of the $f_k$ is assumed.  This is all that
\cite[Lemma~8.4 and Corollary~8.5]{CQ} use: the composite of consecutive
maps is $0$, the sizes are $p+q$ and $q+s$, and the grade conditions on the
factors; the proportionality of \cite[Lemma~8.2]{CQ} needs no stability.
Exactness enters only in Section~\ref{sec:ideal}, where
\cite[Theorem~6.8]{CQ} produces the ranks, stability and the grade
conditions from it.  By Proposition~\ref{prop:dk} the grade conditions are
statements about the $f_k$, and in coordinates they are the hypotheses
of~\cite{CQ} for any bases.

\begin{proof}
\emph{Uniqueness.}  If $a_{k+1}$ is regular, \eqref{eq:star} determines
$a_k$.  This is a proposition, so we may choose bases; let $\rho_k$ be
the induced basis of the line $R_k$ and $\beta=a_{k+1}(\rho_{k+1})$, in
coordinates.  The element $w_x:\rho\mapsto x\wedge a_{k+1}(\rho)$ of $R_k$
is $\mathrm{top}(x\wedge\beta)\,\rho_k$.  If $a_k$ and $a'_k$ both
satisfy~\eqref{eq:star}, then $\mathrm{top}(x\wedge\beta)\,(a_k(\rho_k)-
a'_k(\rho_k))=0$ for all $x$.  Taking for $x$ the basis vectors $e_I$ of
$\Ex^{r_k}R^{s_k}$, the scalars $\mathrm{top}(e_I\wedge\beta)$ are, up to
sign, the coordinates of $\beta$, which form a regular family because
$a_{k+1}$ is regular.  So every coordinate of $a_k(\rho_k)-a'_k(\rho_k)$
is $0$, and $a_k=a'_k$ since $\rho_k$ is a basis.  By descending induction
from $a_{n+1}=\mathrm{id}$, a structure is unique.  Being a structure is
therefore a proposition, and by unique choice existence may be proved
after choosing bases.

\emph{Existence.}  Choose bases of the $F_k$.  They give bases of the
lines $R_k$.  By \cite[Corollary~8.5]{CQ} there are $\beta_k$ with
$\Ex^{r_k}(A_k)=\beta_k\hat\beta_{k+1}^T$, normalised as in
Proposition~\ref{prop:dictionary} below.  Define $a_k$ by
$a_{k}(\text{basis element})=\beta_k$.  By
Proposition~\ref{prop:dictionary}, \eqref{eq:star} holds, and $a_k$ is
regular since $\beta_k$ is ($\Gr(\beta_k)\geq2$ for $k\geq2$
in~\cite{CQ}).
\end{proof}

The comparison with~\cite{CQ} is a computation with the canonical
pairing.  For $x\in\Ex^qR^N$ and $y\in\Ex^{p}R^N$, $p+q=N$, let
$\mathrm{top}(x\wedge y)$ be the coefficient of $e_1\wedge\dots\wedge
e_N$.  Define $\hat y\in\Ex^qR^N$ by $\hat y_I=\epsilon(I,I^c)\,y_{I^c}$
for $|I|=q$, where $e_I\wedge e_{I^c}=\epsilon(I,I^c)\,e_1\wedge\dots\wedge
e_N$.  Then $\mathrm{top}(x\wedge y)=\langle x,\hat y\rangle$, the
coordinatewise scalar product.  For a Hodge dual $y^*$ defined with
$\epsilon(I^c,I)$ in place of $\epsilon(I,I^c)$ one has
$\hat y=(-1)^{pq}y^*$; this is the sign that the factorisation
of~\cite[Lemma~8.4]{CQ} depends on.

\begin{proposition}[dictionary]\label{prop:dictionary}
Let $E\xrightarrow{A}R^N\xrightarrow{B}G$ in coordinates, $\beta\in\Ex^p
R^N$, and $\gamma\in\Ex^q G$, $N=p+q$.  The condition
$\Ex^q(B)(x)=\mathrm{top}(x\wedge\beta)\,\gamma$ for all $x$ of degree
$q$ (that is, \eqref{eq:star} in coordinates) holds if and only if
$\Ex^q(B)=\gamma\hat\beta^T$.  So \cite[Lemma~8.4]{CQ}, every stage of
\cite[Corollary~8.5]{CQ}, and the Cayley determinant are instances
of~\eqref{eq:star}, with $\hat\beta$ in place of $\beta^*$.
\end{proposition}

Thus no sign appears: \eqref{eq:star} never compares $\hat\beta$ with the
Hodge dual.  The chain of~\cite[Corollary~8.5]{CQ}, normalised by
$\beta_{n+1}=1$ and $\alpha_{k}=\hat\beta_{k}$, is itself unique: the
type of normalised chains is contractible, by the same uniqueness
argument in coordinates (a factor $\gamma$ in
$\Ex^q(B)=\gamma\hat\beta^T$ is unique when $\hat\beta$ is regular).

\begin{definition}
$\Det(F):=R_1$, a line, and $a_{\Det}:=a_1:\Det(F)\to\Ex^{r_1}F_0$, for
the unique structure.
\end{definition}

\begin{theorem}\label{thm:det}
Under the hypotheses of Theorem~\ref{thm:structure}, assume moreover that
$F_0$ has rank $r_1$.  Then $\Ex^{r_1}F_0$ is a line, $a_{\Det}$ is
injective, and for any bases, $a_{\Det}$ at the basis element of
$\Det(F)$ determined by the bases is $\delta\,\omega$, where $\omega$ is
the basis element of $\Ex^{r_1}F_0$ and $\delta$ is the Cayley
determinant of~\cite{CQ}.  The image of $a_{\Det}$, i.e.\ the ideal of
values $\lambda(a_{\Det}(\varphi))$ for $\lambda$ a linear form, is the
ideal $(\delta)$; it is the gcd ideal of $D_{r_1}(f_1)$: the elements
divisible by some greatest common divisor.
\end{theorem}

\begin{proof}
In coordinates the structure is the one built from the chain
of~\cite[Corollary~8.5]{CQ}, by uniqueness; its last map sends the basis
element to $\delta\omega$.  Every $\varphi$ is $c$ times the basis
element, so the image is $(\delta)$, and $\delta$ is a regular greatest
common divisor of the $r_1$-minors~\cite[Theorem~8.7]{CQ}.  Injectivity:
$\mathrm{top}(a_{\Det}(c\,\cdot\,\text{basis}))=c\delta$ and $\delta$
is regular.  The gcd ideal is defined without choosing a gcd, and
equals $(g)$ for every gcd $g$.
\end{proof}

\section{Resolutions of an ideal and the greatest common divisor}\label{sec:ideal}

A \emph{finite free resolution} of the ideal $(a_1,\dots,a_N)$ is an
exact complex of finite free modules
\[0\to F_n\xrightarrow{f_n}\cdots\xrightarrow{f_2}F_1\xrightarrow{f_1}F_0\]
($f_n$ injective, exact at each $F_k$, $1\le k<n$) with a linear
bijection $\iota:F_0\cong R$, such that the image of $\iota f_1$ is the
ideal $(a_1,\dots,a_N)$.  The terms are only merely free: no basis is part
of the data.

\begin{proposition}\label{prop:ranks}
Assume $1\neq0$.  The rank data of a resolution, namely the ranks $s_k$
of the $F_k$ and numbers $r_k$ with $s_k=r_k+r_{k+1}$ such that $f_k$ is
stable of rank $r_k$, form a proposition, and are merely given, hence
given.  The hypotheses of Theorem~\ref{thm:structure} hold.  Moreover,
$r_1=1$ unless all $a_j$ are $0$.
\end{proposition}

\begin{proof}
Uniqueness: free ranks and ranks of stable maps are unique since
$1\neq0$.  Existence: after choosing bases,
\cite[Theorem~6.8]{CQ} gives the ranks, stability and
$\Gr(D_{r_k}(f_k))\geq k$; by Proposition~\ref{prop:dk} these hold
without bases.  Since $s_0=1$, $r_1\in\{0,1\}$ or $1=0$: if $r_1\geq2$,
$D_{r_1}(f_1)$ is regular but $\Ex^{2}$ of a rank-one module is $0$.  If
$r_1=0$ then $\Ex^1f_1=0$, so every $1\times1$ minor of $f_1$ is $0$ in
any bases, hence $\iota f_1=0$ and all $a_j$ are $0$.  The case split is
on the number $r_1$, so no choice is involved.
\end{proof}

Let now $(a_j)$ be a nonzero ideal: not all $a_j$ are $0$.  Then $r_1=1$
and Theorem~\ref{thm:det} applies.  The bijection $\iota$, followed by
$R\cong R^1$, is a basis of $F_0$ that comes with the resolution, and
the top coordinate in that basis is an isomorphism
$J:\Ex^1F_0\cong R$.  Define
\[a_1:=J\circ a_{\Det}:\Det(F)\to R.\]
This is the map $a_1:R_1\to\Ex^{1}F_0=R$ of~\cite[Theorem~3.1]{BE74}.

A \emph{greatest common divisor} of $(a_j)$ is $g$ dividing all $a_j$ such
that every common divisor of the $a_j$ merely divides $g$; it is
\emph{regular} if $gx=0$ implies $x=0$.  For regular $g$ the cofactors
are unique, so ``$g$ is a regular gcd'' is a proposition; write
$\RegGcd$ for the set of regular gcds.

\begin{theorem}\label{thm:gcd}
Let $F$ be a finite free resolution of a nonzero ideal $(a_1,\dots,a_N)$,
$1\neq0$.  Then $a_1:\Det(F)\to R$ is injective, its image is the gcd
ideal of $(a_j)$, and
\[\Bases(\Det F)\;\simeq\;\RegGcd,\qquad\varphi\mapsto a_1(\varphi).\]
Both sides are $R^\times$-torsors and the map is equivariant.
\end{theorem}

\begin{proof}
$J$ is a bijection, so $a_1$ is injective by Theorem~\ref{thm:det}, and
its image is the image of $a_{\Det}$, since every linear form on the line
$\Ex^1F_0$ is $J$ times a scalar.  The image of $a_{\Det}$ is the gcd
ideal of the $1\times1$ minors of $f_1$ in a basis, and with the
identification $\iota$ these generate $(a_j)$, so the gcd ideal is that of
$(a_j)$ (a gcd depends only on the ideal).

For any injective linear $a:K\to R$, $\varphi\mapsto a(\varphi)$ is a
bijection between the bases of $K$ and the regular generators of the
image of $a$.  If $\varphi$ is a basis, $a(\varphi)x=a(x\varphi)=0$ gives
$x\varphi=0$, so $x=0$; and $a(c\varphi)=a(\varphi)c$.  Conversely a
regular generator $g$ is $a(\varphi)$ for a unique $\varphi$ ($a$
injective, unique choice); $c\varphi=c'\varphi$ gives $gc=gc'$, so $c=c'$;
and $a(\psi)=gt=a(t\varphi)$ gives $\psi=t\varphi$.

It remains that the regular generators of the gcd ideal are the regular
gcds.  A gcd $g$ generates it.  Conversely let $g$ be a regular generator
and $g_0$ a gcd (merely, by Theorem~\ref{thm:det}).  Then $g_0=gt$ and
$g=g_0u$; so $g$ divides the $a_j$, and every common divisor divides
$g_0$, hence $g$.  Both sides are propositions in $g$, so the two
subtypes of $R$ are equal.  $R^\times$ acts freely and transitively on
regular gcds ($g'=gu$, $g=g'v$ give $vu=1$ since $g$ is regular), and
$a_1(\varphi u)=a_1(\varphi)u$.
\end{proof}

Constructively one cannot split into ``all $a_j$ are $0$'' and ``not all
$a_j$ are $0$''.  A resolution supplies the case distinction: the number
$r_1\in\{0,1\}$ of Proposition~\ref{prop:ranks} is computed from the
resolution.

\begin{proposition}\label{prop:decide}
Let $F$ be a finite free resolution of $(a_1,\dots,a_N)$, $1\neq0$.
Then either all $a_j$ are $0$ (when $r_1=0$), or the regular greatest
common divisors of the $a_j$ form an $R^\times$-torsor, namely the values
of $a_1$ at the bases of $\Det(F)$ (when $r_1=1$).
\end{proposition}

\begin{corollary}
If $1\neq0$ and every ideal $(a,b)$ has a finite free resolution, then for
all $a,b$: either $a=b=0$, or the regular gcds of $a,b$ form an
$R^\times$-torsor.  In both cases $a,b$ merely have a gcd, and the gcd
ideal is merely principal.
\end{corollary}

This is \cite[Corollary~8.10]{CQ} with the gcd as a canonical object.  For
the zero ideal the gcd is $0$, which is not regular unless $1=0$, so
Theorem~8.9 of~\cite{CQ} is to be read with this case distinction.

\section{Independence of the resolution}

\begin{lemma}\label{lem:compare}
Let $a:K\to R$ and $a':K'\to R$ be injective linear maps with the same
image.  There is a unique linear isomorphism $e:K\cong K'$ with $a'e=a$.
\end{lemma}

\begin{proof}
For $x\in K$, $a(x)$ is in the image of $a'$, so merely $a(x)=a'(y)$;
such a $y$ is unique since $a'$ is injective, so it is given by unique
choice, and we put $e(x)=y$.  Linearity: $a'(e(x+x'))=a(x)+a(x')=
a'(e(x)+e(x'))$, and $a'$ is injective; similarly for scalars.  The
inverse is built in the same way from $a'$ and $a$.  If also $a'e'=a$,
then $a'(e'(x))=a'(e(x))$, so $e'=e$.
\end{proof}

\begin{theorem}\label{thm:compare}
Let $1\neq0$ and let $F,G$ be finite free resolutions, of any lengths, of
the same nonzero ideal $(a_1,\dots,a_N)$.  There is a unique isomorphism
$e_{F,G}:\Det(F)\cong\Det(G)$ with $a_G\circ e_{F,G}=a_F$, where $a_F$,
$a_G$ are the maps $a_1$ of $F$ and $G$.  Moreover $e_{F,F}=\mathrm{id}$
and $e_{G,H}\circ e_{F,G}=e_{F,H}$.
\end{theorem}

\begin{proof}
By Theorem~\ref{thm:gcd}, $a_F$ and $a_G$ are injective with image the gcd
ideal of $(a_j)$, which depends only on the $a_j$.  Lemma~\ref{lem:compare}
gives $e_{F,G}$.  The identity and $e_{G,H}e_{F,G}$ satisfy the defining
equations of $e_{F,F}$ and $e_{F,H}$, so they are equal to them by
uniqueness.
\end{proof}

So the line with its injection depends canonically on the ideal, not on
the resolution.

Fix $N$, $a$ and the length $n$.  An \emph{isomorphism} between two
resolutions $(F,f,\iota)$ and $(F',f',\iota')$ of $(a_j)$ is a family of
linear bijections $\varphi_k:F_k\cong F'_k$ with
$\varphi_{k-1}f_k=f'_k\varphi_k$ and $\iota'\varphi_0=\iota$.

\begin{theorem}\label{thm:natural}
An isomorphism of resolutions $F\cong F'$ is an equality $F=F'$.
Consequently, for a nonzero ideal, transport along it gives
$e:\Det(F)\cong\Det(F')$ with $a'_1\circ e=a_1$, and $e=e_{F,F'}$.
\end{theorem}

\begin{proof}
By induction on isomorphisms (the type of families
$\sum_{F'}\prod_k(F_k\cong F'_k)$ is contractible) we may assume
$F'_k=F_k$ and $\varphi_k=\mathrm{id}$.  Then $f'_k=f_k$ and
$\iota'=\iota$ pointwise, hence as linear maps, and the remaining data
(injectivity, exactness, the image condition, the mere existence of
bases) are propositions.  For an equality $F=F'$ the second statement
holds by transport from $e=\mathrm{id}$, and $e=e_{F,F'}$ by the
uniqueness in Theorem~\ref{thm:compare}.
\end{proof}

Every construction on resolutions is natural in the same way.  The same
argument for rings gives, for instance, the invariance of the true grade
under ring isomorphisms as a single transport.

\section{Wiebe's theorem without bases}

Let $V:F\to F$ be an endomorphism of a finite free module of rank $n$.
$\Ex^nV$ acts on the line $\Ex^nF$; the scalar $\Delta$ with
$\Ex^nV(u)=\Delta u$ for all $u$ is unique (the line has a basis
$u_0$, and $\Delta=J(\Delta u_0)$ for an isomorphism $J$ with
$J(u_0)=1$), and exists: in a basis with matrix $A$,
$\Ex(A)(\omega)=\det(A)\,\omega$.  So $\Delta=\det V$ is defined without
bases and is $\det A$ in every basis.

For linear forms $\alpha,\beta:F\to R$ and a finitely generated ideal
$J_0$, write $(J_0,\beta)$ for the ideal of $x$ merely congruent to some
$\beta(v)$ modulo $J_0$, and similarly $(J_0,\beta,\alpha)$ and
$(J_0,\beta,\Delta)$.

\begin{theorem}\label{thm:wiebe}
Assume $\beta\equiv\alpha\circ V$ modulo $J_0$ and $\Gr(\beta)\geq n$
modulo $J_0$ (for the coefficients of $\beta$ in one, equivalently every,
basis).  Then modulo $(J_0,\beta)$: if $x\Delta\in(J_0,\beta)$ then
$x\in(J_0,\beta,\alpha)$, and if $x\alpha(w)\in(J_0,\beta)$ for all $w$
then $x\in(J_0,\beta,\Delta)$.
\end{theorem}

\begin{proof}
The conclusions are propositions; choose a basis.  A linear form
$\lambda$ is $v\mapsto\sum_kc_k(ev)_k$ with $c_k=\lambda(e^{-1}\delta_k)$,
so the image of $\lambda$ is the ideal $(c)$.  With $a,b$ the coefficient
vectors of $\alpha,\beta$ and $A$ the matrix of $eVe^{-1}$ (column $j$
holds the coordinates of the image of $e_j$),
$b_j=\beta(e^{-1}e_j)\equiv\alpha(Ve^{-1}e_j)=\sum_ia_iA_{ij}$ modulo
$J_0$, and $\Delta=\det A$.  This is the situation
of~\cite[Theorem~11]{CT} relative to $J_0$, with $V=A$, which gives the
conclusions for the lists $(J_0,b)$, $(J_0,b,a)$, $(J_0,b,\Delta)$; these
are the ideals of the forms.
\end{proof}

The statement involves no basis, so it is invariant under any change of
basis; in~\cite{CT} the corresponding invariance under elementary
operations is part of the proof, which is used unchanged here.

\section{What is proved and what is checked}

All statements above are checked in the proof assistant Unit
(\url{https://github.com/coquand/unit}), designed by T.~Coquand and
implemented by Claude (Anthropic), in the directory \code{sag/cffru}.  The
matrix-level results of~\cite{CQ} and~\cite{CT} are formalised in
\code{sag/cffr}, \code{sag/koszul} and \code{sag/wiebe}, and the
structure identity principle in \code{sag/sip}; they are used unchanged.
The whole tree checks in a few seconds.

\medskip
{\small\begin{tabular}{ll}
statement & file\\\hline
structure identity (generic; modules, rings) & \code{sag/sip}: \code{SipGeneric}, \code{ModuleSIP2}; \code{RingSIP}\\
$\Ex M$, $\Ex^kM$, the top line, $\Hom$ of lines & \code{ExteriorFree}, \code{TopLine}, \code{Lines}\\
Proposition~\ref{prop:dk}, ranks & \code{DkFree}, \code{Ranks}\\
Proposition~\ref{prop:dictionary} & \code{Dictionary}\\
Theorem~\ref{thm:structure} & \code{StageStar}, \code{Structure}\\
Theorem~\ref{thm:det} & \code{DetLine}\\
Proposition~\ref{prop:ranks} & \code{DetFFR}\\
Theorem~\ref{thm:gcd}, torsors & \code{GcdTorsor}, \code{DetGcd}\\
Proposition~\ref{prop:decide}, Lemma~\ref{lem:compare}, Theorem~\ref{thm:compare} & \code{DetCompare}\\
Theorem~\ref{thm:natural} & \code{NaturalFFR}, \code{DetCompare}\\
contractible chains of~\cite[Cor.~8.5]{CQ} & \code{Factorisations}\\
$\det$ of an endomorphism, Theorem~\ref{thm:wiebe} & \code{DetEnd}, \code{WiebeFree}
\end{tabular}}

\medskip
Not treated: part (b) of~\cite[Theorem~3.1]{BE74}
($\mathrm{Rad}\,I(a_k)=\mathrm{Rad}\,I(f_k)$), whose proof uses prime
ideals; and a proof of Wiebe's theorem that does not pass through
coordinates.

\begin{thebibliography}{9}
\bibitem{BE73} D.~A.~Buchsbaum, D.~Eisenbud, \emph{What makes a complex
  exact?}, J.~Algebra 25 (1973), 259--268.
\bibitem{BE74} D.~A.~Buchsbaum, D.~Eisenbud, \emph{Some structure
  theorems for finite free resolutions}, Adv.~Math.~12 (1974), 84--139.
\bibitem{CL} T.~Coquand, H.~Lombardi, \emph{R\'esolutions libres finies.
  M\'ethodes constructives}, Calvage \& Mounet, 2024.
\bibitem{CQ} T.~Coquand, C.~Quitt\'e, \emph{Constructive finite free
  resolutions}, manuscripta math.~137 (2012), 331--345.
\bibitem{CT} T.~Coquand, C.~T\^ete, \emph{An elementary proof of Wiebe's
  theorem}, J.~Algebra 499 (2018).
\bibitem{E} M.~H.~Escard\'o, \emph{Introduction to Univalent Foundations
  of Mathematics with Agda}, lecture notes, 2019.
\bibitem{R} E.~Rijke, \emph{Introduction to Homotopy Type Theory},
  Cambridge University Press, 2025.
\bibitem{HoTT} The Univalent Foundations Program, \emph{Homotopy Type
  Theory: Univalent Foundations of Mathematics}, Institute for Advanced
  Study, 2013.
\bibitem{N} D.~G.~Northcott, \emph{Finite Free Resolutions}, Cambridge
  University Press, 1976.
\end{thebibliography}
\end{document}
